\documentclass[11pt,a4paper]{article}

\usepackage[margin=1.8cm]{geometry}
\usepackage{enumitem}
\usepackage[hidelinks]{hyperref}
\usepackage[T1]{fontenc}
\usepackage{titlesec}

\setlength{\parindent}{0pt}
\pagenumbering{gobble}
\titleformat{\section}{\large\bfseries}{}{0em}{}[\tit                                                                                                                                      lerule]

\begin{document}

%==================== HEADER ====================
\begin{center}
{\Large \textbf{ING. MIGUEL ANGEL ANAY GOMEZ}}\\
Lima, Perú\\
\textbf{Software Architect | Microservices | IAM | API Security | AWS | .NET | DevSecOps}\\

\href{mailto:contacto@miguel-anay.nom.pe}{contacto@miguel-anay.nom.pe} \;|\;
\href{https://www.linkedin.com/in/miguel-anay/}{LinkedIn} \;|\;
+51 936 327 402 \;|\;
\href{https://portfolio.miguel-anay.nom.pe}{Portafolio}
\end{center}

\vspace{0.3cm}

%==================== PERFIL ====================
\section*{Perfil Profesional}

Arquitecto de Software con más de \textbf{12 años de experiencia} en diseño e implementación de soluciones empresariales distribuidas.  
Especializado en \textbf{arquitecturas de microservicios, seguridad en APIs, modelos de autenticación y autorización, e identidad de máquina (Machine Identity)}.  

Experiencia diseñando arquitecturas end-to-end, definiendo estándares técnicos, esquemas de integración y lineamientos de seguridad bajo principios de \textbf{Zero Trust}.  
Participación activa en definición de modelos IAM, autenticación basada en tokens (OAuth2 / OIDC / JWT) y autenticación servicio-a-servicio mediante \textbf{mTLS}.  


%==================== EXPERIENCIA ====================
\section*{Experiencia Profesional}

\vspace{6pt}

\textbf{Yiwu Import Corporation} \hfill \textit{Abr 2025 – Dic 2025}\\
\textit{Arquitecto de Tecnología / Líder Técnico}

\begin{itemize}[leftmargin=*]
  \item Diseño de arquitectura basada en microservicios desplegada en AWS (EC2 / Kubernetes).
  \item Implementación de autenticación basada en JWT y OAuth2 para APIs internas y externas.
  \item Definición de modelo de autorización RBAC para servicios distribuidos.
  \item Implementación de comunicación servicio-a-servicio mediante mTLS.
  \item Diseño de lineamientos de seguridad para APIs y control de acceso.
  \item Contenerización con Docker y definición de estándares de despliegue CI/CD.
\end{itemize}

\textbf{Logros:}
\begin{itemize}[leftmargin=*]
    \item Diseño e implementación de una arquitectura de microservicios segura en AWS, incorporando autenticación basada en tokens, comunicación servicio-a-servicio cifrada y estandarización DevSecOps, asegurando escalabilidad y reducción de riesgos operativos.

\end{itemize}

\vspace{6pt}

\textbf{Encora – Cliente Profuturo AFP} \hfill \textit{Nov 2020 – Abr 2025}\\
\textit{Technical Lead / Arquitecto de Software}

\begin{itemize}[leftmargin=*]
  \item Diseño de arquitectura técnica para plataformas digitales basadas en microservicios.
  \item Definición de estándares de autenticación y autorización para servicios internos.
  \item Implementación de flujos OAuth2 y OpenID Connect para gestión segura de identidad.
  \item Diseño de APIs seguras con validación descentralizada de JWT.
  \item Coordinación con equipos de seguridad para cumplimiento de lineamientos IAM.
  \item Elaboración de diagramas de arquitectura, integración y flujos de identidad.
\end{itemize}


\textbf{Logros:}
\begin{itemize}[leftmargin=*]
    \item Definición e implementación de lineamientos de arquitectura y seguridad para plataformas digitales críticas, fortaleciendo los controles de autenticación y autorización y reduciendo incidencias relacionadas a acceso en más del 40\%.

\end{itemize}

\vspace{6pt}

\textbf{Prosegur} \hfill \textit{Abr 2018 – Mar 2020}\\
\textit{Líder Técnico de Soluciones}

\begin{itemize}[leftmargin=*]
  \item Diseño y desarrollo de \textbf{Web Audit}, plataforma de auditoría de abastecimiento de ATM con control de acceso basado en roles (RBAC).
  \item Implementación de mecanismos de autenticación segura para usuarios internos mediante validación de credenciales y tokens.
  \item Implementación de trazabilidad y auditoría de accesos a información sensible.
\end{itemize}

\textbf{Logros:}
\begin{itemize}[leftmargin=*]
  \item Fortalecimiento de controles de acceso a información sensible operativa.
  \item Reducción de riesgos mediante validaciones y segmentación de permisos por rol.
\end{itemize}


\vspace{6pt}

\textbf{Medios Industriales} \hfill \textit{Abr 2014 – Mar 2018}\\
\textit{Arquitecto de Software .NET}

\begin{itemize}[leftmargin=*]
  \item Desarrollo de sistema \textbf{CTO} para control de tiempos operativos.
  \item Implementación de plataforma \textbf{Balizas} para monitoreo GPS en tiempo real.
\end{itemize}

\textbf{Logros:}
\begin{itemize}[leftmargin=*]
  \item Optimización de recursos y mejora de productividad operativa.
\end{itemize}

%==================== COMPETENCIAS ====================
\section*{Competencias Técnicas}

\textbf{Arquitectura:} Microservicios, Arquitectura Distribuida, Diseño End-to-End, Seguridad en APIs, Zero Trust\\
\textbf{Seguridad:} OAuth2, OpenID Connect, JWT, RBAC, ABAC, mTLS, Machine Identity, SPIFFE\\
\textbf{Cloud:} AWS (EC2, S3, IAM, EKS), Azure\\


%==================== EDUCACION ====================
\section*{Educación}

\textbf{Universidad Nacional de Ingeniería (UNI)}\\
Ingeniero Industrial \hfill \textit{2004 – 2012}\\
CIP: 301084

\vspace{4pt}

\textbf{Universidad Privada del Norte}\\
Diplomado en Gerencia de TI – Disrupción Tecnológica\\
\textit{Ene 2025 – Sep 2025}

%==================== ESPECIALIZACION ====================
\section*{Especialización}

\begin{itemize}[leftmargin=*]
  \item Advanced Data Engineer – DMC Instituto (2025 – Actual)
  \item Gen AI Training Path – Coursera (2025)
  \item Azure Fundamentals AZ-900 – Microsoft (2024)
  \item Arquitectura Hexagonal C\#.NET (2025)
  \item Administración SQL Server
\end{itemize}

%==================== IDIOMAS ====================
\section*{Idiomas}

Español: Nativo\\
Inglés: Intermedio

\end{document}
